\section{Consumption economics：把资源影响反馈给开发者}

云和 AI 都在走向 consumption-based pricing。问题不只是“按量收费是否便宜”，而是计量能否成为开发者的 \term{fitness function}：某次 query、render、model call 或 data transfer 消耗了什么资源，是否影响其他 workload，开发者能否据此优化。

\subsection{Honest metering 与可归因资源}

若计费维度为 $u_i$，单价为 $p_i$，账单为
\[
B=\sum_i u_i p_i.
\]
公式本身不难；困难在于 $u_i$ 是否与开发者操作可对应。一个昂贵 join 若只表现为“数据库整体变慢、账单整体上升”，反馈很差；若系统能指出该 operation 的读取、CPU、内存和延迟影响，开发者才能决定加索引、缓存、拆查询或接受成本。

\lecturefigure{09-consumption-feedback.png}{消费计量只有进入开发决策闭环时才产生工程价值。}{本地字幕 20:00--22:40；概念重绘。}

\begin{knowledgebox}{读图：价格信号必须和质量信号一起看}
只显示费用会诱导过度优化；只显示延迟会隐藏资源浪费。理想界面同时展示 request volume、latency、error、cache hit、CPU/memory、transfer 和 spend，并关联 deployment version。开发者优化的是单位业务结果的成本，而不是单个指标最小化。
\end{knowledgebox}

\subsection{Noisy workload 与隔离}

课堂用数据库作类比，真正要表达的是 workload attribution：共享系统中某个昂贵操作可能拖慢邻居，却没有清晰计量。讲义不把 PostgreSQL 或 DynamoDB 绝对分类为“坏/好系统”，而是比较具体配置能否提供 workload isolation、horizontal scaling 和 operation-level evidence。

\lecturefigure{10-workload-isolation.png}{资源隔离和 operation-level accounting 让性能与成本可解释。}{本地字幕 20:30--22:35；概念重绘。}

\begin{warningbox}{消费计费可能把平台低效转嫁给用户}
按量收费不天然公平。如果平台 runtime 浪费 CPU、cache 策略低效或计量维度不可解释，用户会为平台问题买单。供应商需要公开计量语义、提供优化建议，并让用户区分业务增长与系统回归。
\end{warningbox}

\subsection{本章小结}

Consumption economics 的核心是 attribution。资源、价格和服务质量必须映射回具体 operation 和 deployment，才能形成工程反馈，而不是月底才出现的不可解释账单。

